%!TEX root = ../main_thesis.tex

\chapter{Analisi sintattica e rappresentazione intermedia}
\label{cap:parsing}

Questo capitolo apre la parte dedicata al \emph{front-end} del compilatore, che
comprende le prime fasi della traduzione: il riconoscimento della struttura del
testo sorgente e la sua codifica in una rappresentazione interna su cui le fasi
successive potranno lavorare.

\section{Concetti teorici}
\label{sec:parsing-teoria}

\paragraph{Grammatiche e analisi sintattica.}
La sintassi di un linguaggio è descritta da una \emph{grammatica}: un insieme di
regole che stabiliscono quali sequenze di simboli costituiscono frasi valide.
L'\emph{analisi sintattica} (\emph{parsing}) verifica che il testo sorgente
rispetti tali regole e ne costruisce una rappresentazione strutturata. La
formalizzazione più diffusa è quella delle grammatiche \emph{context-free}
(CFG), spesso scritte in notazione BNF, in cui ogni regola può prevedere più
alternative. Le CFG possono però essere \emph{ambigue}, cioè ammettere più
alberi di derivazione per la stessa frase, il che complica la costruzione del
parser.

\paragraph{Parsing Expression Grammar.}
Le \emph{Parsing Expression Grammar} (PEG)~\cite{ford2004} offrono una
formalizzazione alternativa, pensata direttamente in funzione del
riconoscimento. La differenza fondamentale rispetto alle CFG è la \emph{scelta
ordinata}: quando una regola prevede più alternative, queste vengono provate
nell'ordine in cui sono scritte e viene accettata la prima che ha successo.
Questo rende una PEG intrinsecamente non ambigua. Inoltre una PEG integra
l'analisi lessicale e quella sintattica in un'unica descrizione
(\emph{scannerless}), eliminando la necessità di un analizzatore lessicale
separato: la gestione degli spazi e dei commenti è parte della grammatica
stessa.

\paragraph{La rappresentazione intermedia.}
Il risultato del parsing è una rappresentazione ad
albero del testo sorgente che ne cattura la struttura logica trascurando i dettagli sintattici
superflui. Questa \emph{rappresentazione intermedia} (IR) è la struttura dati su
cui operano le fasi successive del compilatore. Una buona IR rispecchia da
vicino i costrutti del linguaggio, così da rendere naturali le elaborazioni
successive, e conserva le informazioni necessarie alla diagnostica, in
particolare la posizione di ciascun elemento nel sorgente.

La scelta di una PEG per ROOC è motivata proprio da queste proprietà: la non
ambiguità semplifica la definizione della grammatica, e l'approccio
\emph{scannerless} permette di descrivere in un solo file l'intera sintassi del
linguaggio, comprese le sue particolarità come la moltiplicazione implicita.

\section{Applicazione in ROOC: la grammatica PEST}
\label{sec:grammatica-pest}

ROOC definisce la propria sintassi tramite \textbf{PEST}, una libreria Rust che
implementa le PEG. La grammatica è scritta in un file dedicato\footnote{La
grammatica completa è disponibile nel repository del progetto:
\url{https://github.com/Specy/rooc/blob/main/src/parser/grammar.pest}} e viene
tradotta in un parser a tempo di compilazione.

\paragraph{La regola principale.}
La regola di partenza, \texttt{problem}, descrive la struttura complessiva di un
modello: la funzione obiettivo, la sezione dei vincoli e i blocchi opzionali
\texttt{where} e \texttt{define} (Listato~\ref{lst:grammar-problem}).

\begin{lstlisting}[language=pest, caption={La regola principale della grammatica.}, label={lst:grammar-problem}]
problem = {
    SOI ~ nl* ~
    #objective = objective ~ nl+ ~
    (^"s.t." | ^"subject to") ~ nl+ ~
    #constraints = constraint_list ~
    (nl+ ~ ^"where"  ~ #where  = consts_declaration)? ~
    (nl+ ~ ^"define" ~ #define = domains_declaration)? ~
    nl* ~ EOI
}
\end{lstlisting}

\begin{table}[H]
\centering
\begin{tabular}{l p{0.70\textwidth}}
\hline
\textbf{Costrutto} & \textbf{Significato} \\
\hline
\texttt{\#nome = \ldots} & \emph{Tag}: assegna un nome a una sottoparte della regola, rendendola facilmente recuperabile in fase di costruzione dell'IR. \\
\texttt{\textasciitilde} & \emph{Sequenza}: gli elementi devono comparire in successione, nell'ordine indicato. \\
\texttt{\textbar} & \emph{Scelta ordinata}: le alternative sono provate da sinistra a destra e si accetta la prima che ha successo. \\
\texttt{?}, \texttt{*}, \texttt{+} & \emph{Quantificatori}: l'elemento compare rispettivamente zero o una volta, zero o più volte, una o più volte. \\
\texttt{nl} & Regola ausiliaria che riconosce un'interruzione di riga (\emph{newline}). \\
\texttt{\textasciicircum\textquotedbl\ldots\textquotedbl} & Stringa letterale \emph{case-insensitive}: \texttt{\textasciicircum\textquotedbl s.t.\textquotedbl}. \\
\texttt{@\{\ldots\}} & Regola \emph{atomica}: trattata come un unico token terminale. \\
\texttt{\_\{\ldots\}} & Regola \emph{silenziosa}: non compare nell'albero risultante perché serve solo a strutturare la grammatica. \\
\texttt{(\ldots)} & Raggruppamento: applica un operatore (per esempio \texttt{?} o \texttt{*}) a un gruppo di elementi. \\
\texttt{SOI}, \texttt{EOI} & Inizio (\emph{Start Of Input}) e fine (\emph{End Of Input}) dell'input. \\
\hline
\end{tabular}
\caption{Gli operatori della grammatica PEG/PEST utilizzati in ROOC.}
\label{tab:pest-operatori}
\end{table}

La regola \texttt{problem} (Listato~\ref{lst:grammar-problem}) impiega gran parte
degli operatori delle PEG e delle estensioni introdotte da PEST, il cui
significato è riassunto nella Tabella~\ref{tab:pest-operatori}.

\paragraph{Regole atomiche, silenziose e scelta ordinata.}
Il Listato~\ref{lst:grammar-misc} mostra in concreto alcuni di questi operatori:
la regola \texttt{objective\_type} è \emph{atomica} (\texttt{@}) e impiega la
scelta ordinata (\texttt{\textbar}) tra due stringhe \emph{case-insensitive},
mentre \texttt{exp} è \emph{silenziosa} (\texttt{\_}) e combina opzionalità e
ripetizione.

\

\begin{lstlisting}[language=pest, caption={Esempi di regola atomica, silenziosa e di scelta ordinata.}, label={lst:grammar-misc}]
objective_type = @{ ^"min" | ^"max" }
exp = _{ unary_op? ~ exp_leaf ~ (binary_op ~ unary_op? ~ exp_leaf)* }
\end{lstlisting}

\paragraph{Dall'albero di parsing all'IR.}
L'albero prodotto da PEST viene poi visitato e trasformato nella rappresentazione intermedia del modello.

\section{Il \texttt{PreModel} e l'Intermediate Language}
\label{sec:premodel}

La rappresentazione intermedia di ROOC prende il nome di \texttt{PreModel}: è il
modello così come risulta dalla sola analisi sintattica, prima dell'analisi
semantica e della trasformazione. Esso raccoglie le quattro componenti di un
modello, più un riferimento al codice sorgente originale, utile per la
diagnostica.

La funzione obiettivo è rappresentata da un \texttt{PreObjective}, che ne indica
il tipo (minimizzazione o massimizzazione) e l'espressione associata. Ogni
vincolo è un \texttt{PreConstraint}, che contiene le due espressioni messe a
confronto, il tipo di relazione, l'eventuale nome e le eventuali clausole di
iterazione (\texttt{for}) che lo replicano. Sia l'obiettivo sia i vincoli
contengono dunque delle \emph{espressioni}, che costituiscono il cuore dell'IR.

\paragraph{Le espressioni: \texttt{PreExp}.}
Le espressioni sono rappresentate dal tipo enumerativo \texttt{PreExp}, le cui
varianti rispecchiano da vicino i costrutti del linguaggio descritti nel
Capitolo~\ref{cap:linguaggio}: valori primitivi, variabili semplici e composte,
accessi ad array, chiamate di funzione, \emph{block function} (con e senza
\emph{scope}) e operazioni binarie e unarie (Listato~\ref{lst:preexp}).

\

\begin{lstlisting}[language=rust, caption={Versione semplificata dell'enumerazione \texttt{PreExp}.}, label={lst:preexp}]
enum PreExp {
    Primitive(Spanned<Primitive>),
    Variable(Spanned<String>),
    CompoundVariable(Spanned<CompoundVariable>),
    ArrayAccess(Spanned<AddressableAccess>),
    FunctionCall(InputSpan, FunctionCall),
    BlockFunction(Spanned<BlockFunction>),
    BlockScopedFunction(Spanned<BlockScopedFunction>),
    BinaryOperation(Spanned<BinOp>, Box<PreExp>, Box<PreExp>),
    UnaryOperation(Spanned<UnOp>, Box<PreExp>),
    // ...
}
\end{lstlisting}

Essendo un albero, una \texttt{PreExp} può contenere altre \texttt{PreExp}: per
esempio un'operazione binaria racchiude i due sotto-espressioni operandi. Le
variabili composte, come \texttt{x\_i}, sono rappresentate da un
\texttt{CompoundVariable} che ne memorizza il nome e l'elenco delle espressioni
usate come indice; analogamente, l'accesso a un array è descritto da un
\texttt{AddressableAccess}.

\section{Gestione e tracciamento degli errori}
\label{sec:errori-parsing}

Un aspetto trasversale a tutte le fasi del compilatore è la capacità di
segnalare gli errori indicando con precisione il punto del sorgente a cui si
riferiscono.

A questo scopo, ogni elemento della rappresentazione intermedia è annotato con
la propria posizione nel sorgente tramite un contenitore generico,
\texttt{Spanned}, che incapsula un valore insieme alle informazioni sulla sua
posizione (riga, colonna e lunghezza del frammento di testo). Come si è visto,
gran parte delle varianti di \texttt{PreExp} avvolge i propri dati in uno
\texttt{Spanned}, così che ogni nodo dell'albero conservi il punto del sorgente
da cui proviene.

Grazie a questa annotazione, gli errori sono \emph{posizionali}: ogni
segnalazione è collegata al frammento di sorgente che l'ha generata, e può così
essere presentata indicando dove l'errore ha origine. La stessa infrastruttura
permette di tracciare non solo gli errori di parsing, ma anche quelli delle fasi
successive, come gli errori di tipo e di trasformazione descritti nel
Capitolo~\ref{cap:semantica}.
